% All five plots show the same coordinate window: [-3,3] x [-6,6].
\newcommand{\PolynomialComponentPlot}[3]{%
  \begin{scope}[xshift=#1cm,yshift=-2.2cm,x=0.46cm,y=0.24cm]
    \draw[step=1,black!7,line width=0.25pt] (-3,-6) grid (3,6);
    \draw[->,black!55] (-3.2,0)--(3.3,0) node[right=2pt] {$x$};
    \draw[->,black!55] (0,-6.25)--(0,6.3) node[above=2pt] {$y$};
    \foreach \x in {-3,-2,-1,1,2,3} {
      \draw[black!45] ([yshift=-1pt]\x,0)--([yshift=1pt]\x,0);
    }
    \foreach \x in {-3,3} {
      \node[below=3pt,inner sep=0pt] at (\x,0) {$\x$};
    }
    \foreach \y in {-6,-4,-2,2,4,6} {
      \draw[black!45] ([xshift=-1pt]0,\y)--([xshift=1pt]0,\y);
    }
    \foreach \y in {-6,6} {
      \node[left=3pt,inner sep=0pt] at (0,\y) {$\y$};
    }
    \begin{scope}
      \clip (-3,-6) rectangle (3,6);
      \draw[#2,line width=1.1pt,domain=-3:3,samples=121]
        plot (\x,{#3});
    \end{scope}
  \end{scope}%
}

\begin{tikzpicture}[>=Latex,line cap=round]
  % Each term is centered immediately above its own graph.
  \node[anchor=west,inner sep=0pt] at (-8.5,0.35) {$p(x)=$};
  \node[text=green!45!black,inner sep=1pt] at (-6,0.35) {$x^3$};
  \node[inner sep=1pt] at (-4,0.35) {$+$};
  \node[text=black!75,inner sep=1pt] at (-2,0.35) {$(-x^2)$};
  \node[inner sep=1pt] at (0,0.35) {$+$};
  \node[text=red!75!black,inner sep=1pt] at (2,0.35) {$(-x)$};
  \node[inner sep=1pt] at (4,0.35) {$+$};
  \node[text=orange!85!black,inner sep=1pt] at (6,0.35) {$1$};

  % Parenthesize the substituted variable before exponentiation: otherwise
  % a negative sample such as -2^2 is parsed as -(2^2), not (-2)^2.
  \PolynomialComponentPlot{-6}{green!45!black}{(\x)^3}
  \PolynomialComponentPlot{-2}{black!75}{-(\x)^2}
  \PolynomialComponentPlot{2}{red!75!black}{-(\x)}
  \PolynomialComponentPlot{6}{orange!85!black}{1}

  % Four additive contributions flow down into the combined graph. The plus
  % signs sit beside the arrows, with transparent backgrounds and clear space.
  \draw[->,green!45!black,line width=0.8pt]
    (-6,-4.05) to[out=-90,in=90]
    node[pos=0.48,above=6pt,font=\Large,inner sep=0pt] {$+$} (-3,-5.65);
  \draw[->,black!75,line width=0.8pt]
    (-2,-4.05) to[out=-90,in=90]
    node[pos=0.48,left=5pt,font=\Large,inner sep=0pt] {$+$} (-1,-5.65);
  \draw[->,red!75!black,line width=0.8pt]
    (2,-4.05) to[out=-90,in=90]
    node[pos=0.48,right=5pt,font=\Large,inner sep=0pt] {$+$} (1,-5.65);
  \draw[->,orange!85!black,line width=0.8pt]
    (6,-4.05) to[out=-90,in=90]
    node[pos=0.48,above=6pt,font=\Large,inner sep=0pt] {$+$} (3,-5.65);

  \begin{scope}[yshift=-8.4cm,x=1.15cm,y=0.42cm]
    \draw[step=1,black!7,line width=0.25pt] (-3,-6) grid (3,6);
    \draw[->,black!55] (-3.2,0)--(3.3,0) node[right=2pt] {$x$};
    \draw[->,black!55] (0,-6.25)--(0,6.3) node[above=2pt] {$y$};
    \foreach \x in {-3,-2,-1,1,2,3} {
      \draw[black!45] ([yshift=-1.5pt]\x,0)--([yshift=1.5pt]\x,0);
    }
    \foreach \x in {-3,-2,1,2,3} {
      \node[below=3pt,inner sep=0pt] at (\x,0) {$\x$};
    }
    \node[below right=3pt,inner sep=0pt] at (-1,0) {$-1$};
    \foreach \y in {-6,-4,-2,2,4,6} {
      \draw[black!45] ([xshift=-1.5pt]0,\y)--([xshift=1.5pt]0,\y);
      \node[left=3pt,inner sep=0pt] at (0,\y) {$\y$};
    }
    \IfClozeTF{}{
      \begin{scope}
        \clip (-3,-6) rectangle (3,6);
        \draw[THMBlue,line width=1.3pt,domain=-3:3,samples=241]
          plot (\x,{(\x)^3-(\x)^2-(\x)+1});
      \end{scope}
    }
  \end{scope}
\end{tikzpicture}
